\section{Closed-Loop Diagnostics}
\label{sec:diagnostics}

Matched continuations distinguish state prediction from prediction of a
control intervention. From one recorded state, 20 versus 30 s of holding
reduces waiting plus onboard exposure by 125 person-s under native
continuation and by 70 person-s under rolling control over the same 582-s
window. Actual cumulative CV$^2$ also decreases in both cases. Total holding
is unchanged within each pair: its location and passenger exposure change.
Reoptimization therefore attenuates this local benefit without reversing it.

Two diagnostics explain why point-forecast improvements need not improve
control. First, a predicted downstream arrival lay beyond the physical
comparison window; removing that event changed the predicted CV$^2$ difference
from approximately $+0.936\times10^{-6}$ to $-0.201\times10^{-6}$. A
cross-window sign disagreement could not establish a wrong action ranking.
Second, improved motion predictions reduced six continuous holding-target
errors below 0.562 s, yet some predictions remained across discrete-action
boundaries. For example, a predicted target of 7.9019 s and actual value of
7.3404 s straddle the 7.5-s boundary. Relevant validation must include the
projected response to an intervention and its actual continuation, not only
absolute timing error.

\section{Discussion and Evidence Boundaries}
\label{sec:discussion}

\subsection{Observation and Configuration Semantics}

A selected queue forecast of 7.0694 passengers versus a retrospective count
of one motivated an audit rather than a fitted correction. Across 441 queries
from 73 other APC intervals, interval-weighted mean error was $-0.326$
passengers and became $-1.570$ for gaps of at least 360 s. No additional eligible
target-stop observation supplied the missing arrivals, so a uniform downward
correction was unsupported. The point-correction branch was stopped.

The budget experiment gives operational meaning to configuration. The
retention floor $\rho_B$ limits the use of vehicle time while preserving
explicit timetable obligations. Its monotonic effect on a same-state
allowance does not ensure a monotonic full-day outcome. The prespecified
settings demonstrate actual service and circulation consequences, including
the regularity cost of tighter retention. They provide a concrete interface
for examining dispatch preferences. Human usability has not been measured.

\subsection{Limitations and Scope of the Claim}
\label{sec:limitations}

The learned controller's advantage over strong manual joint remains unresolved
at both matched retention settings. At retention 0.25, waiting and journey
improve only on the fourth day; intervals spanning zero establish neither
equivalence nor noninferiority. The older cross-budget contrast combines
configuration and upper-policy changes. The matched harmonic ablation supports
a modest journey increment at default retention while changing estimator
capacity and both actor and plan-law inputs. A direct high-frequency fast-control
pathway and complete LF/HF specialization remain unestablished. The fixed-fleet
implementation does not allocate vehicles or drivers by frequency band.

All full-day findings concern reused development days; independent confirmation
remains necessary. The additional manual-budget comparison uses one frozen
learning seed, whereas the default harmonic comparison reports three seeds;
the two studies do not constitute a crossed three-seed budget experiment.
The local continuation observes only seven of 24 resource endpoints and cannot
establish an all-day guarantee.

An enabled observation-vector field described as a queue estimate reads the
simulator's waiting population. It is ineligible as an additional APC
observation. The examined fast native rule does not consume that component;
future learned fast policies require an explicit observation-use audit.

A complete layer comparison should separate fixed versus dynamic timetable
control from reference versus adaptive holding under matched budgets and
observation access. A separately configured and trained TransitDuet reference
is also needed; no manual mode here is relabeled as TransitDuet. Representation-
capacity controls and independent demand days would strengthen attribution
and generalization. The current evidence supports executable coordination,
a modest harmonic increment, and measurable budget-setting tradeoffs.
